\documentclass[sigconf,nonacm]{acmart}

% Remove ACM-specific footer elements for class submission
\settopmatter{printacmref=false}
\renewcommand\footnotetextcopyrightpermission[1]{}
\pagestyle{plain}
\setcopyright{none}

% Packages
\usepackage{booktabs}
\usepackage{graphicx}

\begin{document}

\title{Predicting Wildfire Ignition Cause in Canada\\Using Machine Learning}

\author{[Student Name 1]}
\affiliation{%
  \institution{Technical Expert -- Computer Science}
  \institution{Carleton University}
}
\email{student1@cmail.carleton.ca}

\author{[Student Name 2]}
\affiliation{%
  \institution{Domain Expert -- Geography}
  \institution{Carleton University}
}
\email{student2@cmail.carleton.ca}

\author{[Student Name 3]}
\affiliation{%
  \institution{Domain Expert -- Environmental Science}
  \institution{Carleton University}
}
\email{student3@cmail.carleton.ca}

\begin{abstract}
This project develops a machine learning classification model to predict whether a wildfire in Canada was naturally caused (lightning) or human-caused. Using the Canadian National Fire Database (1972--2024) containing 14,485 fires with labeled causes, combined with Environment Canada climate data, we train and evaluate multiple classifiers. Preliminary analysis achieved 91.6\% accuracy and 0.891 ROC-AUC, demonstrating strong feasibility. Results will provide actionable insights for fire prevention resource allocation and targeted awareness campaigns.
\end{abstract}

\maketitle

\section{Problem Statement}

Wildfires represent one of Canada's most significant natural hazards. The 2023 fire season alone burned over 16 million hectares---more than double any previous record. Understanding wildfire \textit{cause} is critical for prevention: while natural fires (lightning-ignited) are largely unavoidable, human-caused fires are entirely preventable through targeted interventions.

The Canadian National Fire Database (CNFDB) contains records of over 16,000 large fires (>200 hectares) from 1972--2024, with cause labels assigned by provincial fire management agencies. Approximately 86\% of large fires are naturally caused, while 14\% are human-caused. However, the relationship between environmental conditions and fire causation remains underexplored.

\textbf{Research Question:} Can we accurately predict whether a wildfire was naturally or human-caused based on temporal, geographic, and climate features? Which factors most strongly distinguish these fire types?

\section{Motivation}

This research addresses a practical problem with significant societal impact:

\begin{enumerate}
    \item \textbf{Fire Prevention}: Human-caused fires are preventable. Identifying high-risk conditions enables targeted prevention campaigns.
    
    \item \textbf{Resource Allocation}: Fire agencies can optimize deployment by understanding when and where different fire types occur.
    
    \item \textbf{Policy Development}: Evidence-based insights inform regulations on campfires, industrial activities, and land use.
    
    \item \textbf{Climate Context}: As fire seasons lengthen due to climate change, understanding weather-causation relationships becomes critical.
\end{enumerate}

From a data science perspective, this project offers an excellent opportunity to apply the course concepts---data cleaning, exploratory analysis, feature engineering, and machine learning classification---to a real-world problem with reliable ground truth labels, interpretable features, and meaningful societal implications~\cite{flannigan}.

\section{Objectives}

\subsection{Primary Objective}
Develop a machine learning classifier that predicts wildfire ignition cause (Natural vs. Human) with accuracy exceeding 85\%, using the CNFDB and Environment Canada climate data.

\subsection{Secondary Objectives}
\begin{itemize}
    \item Identify key features distinguishing natural from human-caused fires through feature importance analysis
    \item Analyze regional and seasonal patterns in fire causation across provinces
    \item Integrate climate variables (temperature, precipitation, drought) to improve predictions
    \item Provide actionable recommendations for fire management agencies
\end{itemize}

\section{Data Sources}

\subsection{Canadian National Fire Database}
Source: Natural Resources Canada, Canadian Forest Service~\cite{cwfis}

\begin{itemize}
    \item \textbf{Coverage}: 1972--2024, all provinces and territories
    \item \textbf{Records}: 16,394 large fires; \textbf{14,485 with labeled cause}
    \item \textbf{Target Variable}: \texttt{CAUSE} (N=Natural, H=Human)
    \item \textbf{Features}: Date, province, fire size, detection method
    \item \textbf{Label Quality}: Assigned by professional fire investigators; 88\% of records have known cause
\end{itemize}

\subsection{Environment Canada Climate Data}
Source: Government of Canada~\cite{eccc}

\begin{itemize}
    \item \textbf{Coverage}: Daily observations from weather stations nationwide
    \item \textbf{Variables}: Temperature, precipitation, snow, wind speed
    \item \textbf{Integration}: Match fires to nearest stations for climate context
\end{itemize}

\section{Methodology}

\subsection{Data Preparation}
Load CNFDB fire data, filter to records with known cause labels, acquire climate data for fire-prone regions (BC, AB, ON, QC), and perform quality assessment.

\subsection{Feature Engineering}
\begin{itemize}
    \item \textbf{Temporal}: Day of year (seasonality), month, year, weekend indicator
    \item \textbf{Geographic}: Province encoding
    \item \textbf{Climate}: Temperature anomalies, precipitation deficit
    \item \textbf{Fire}: Log-transformed fire size
\end{itemize}

\subsection{Model Development}
Train and compare: Logistic Regression (baseline), Random Forest, and Gradient Boosting. Address class imbalance (86\%/14\%) using class weights.

\subsection{Evaluation}
Stratified cross-validation with accuracy, precision, recall, F1, and ROC-AUC metrics. Feature importance analysis to interpret results.

\section{Preliminary Results}

A proof-of-concept using fire metadata features (no climate data yet) achieved:

\begin{table}[h]
\small
\centering
\caption{Model Performance (Test Set: 2,878 fires)}
\begin{tabular}{lcc}
\toprule
\textbf{Model} & \textbf{Accuracy} & \textbf{ROC-AUC} \\
\midrule
Logistic Regression & 90.8\% & 0.857 \\
Random Forest & 91.8\% & 0.888 \\
\textbf{Gradient Boosting} & \textbf{91.6\%} & \textbf{0.891} \\
\bottomrule
\end{tabular}
\end{table}

\textbf{Key Finding}: Day of year is the strongest predictor (49\% importance). Natural fires peak in July--August (lightning season), while human fires have broader seasonal distribution. Integrating climate data is expected to further improve performance.

\section{Expected Contributions}

\begin{enumerate}
    \item Validated ML model for wildfire cause prediction
    \item Analysis of temporal and regional causation patterns
    \item Recommendations for prevention resource allocation
    \item Reproducible Python code and methodology
\end{enumerate}

\section{Tools \& Technologies}
Python (pandas, scikit-learn, matplotlib), Jupyter notebooks for analysis, Git for version control.

\section{Team Composition}
\begin{itemize}
    \item \textbf{Technical Expert} ([Name 1]): Machine learning model development, data pipeline, statistical analysis
    \item \textbf{Domain Expert} ([Name 2]): Problem framing, ecological interpretation, fire science context
    \item \textbf{Domain Expert} ([Name 3]): Policy implications, literature review, visualization
\end{itemize}

\bibliographystyle{ACM-Reference-Format}
\begin{thebibliography}{3}

\bibitem{cwfis}
Canadian Forest Service. 2025. Canadian National Fire Database. Natural Resources Canada. \texttt{https://cwfis.cfs.nrcan.gc.ca/datamart}

\bibitem{eccc}
Environment and Climate Change Canada. Historical Climate Data. Government of Canada. \texttt{https://dd.weather.gc.ca/}

\bibitem{flannigan}
Flannigan, M.D., Krawchuk, M.A., de Groot, W.J., Wotton, B.M., and Gowman, L.M. 2009. Implications of changing climate for global wildland fire. \textit{International Journal of Wildland Fire}, 18(5), 483--507.

\end{thebibliography}

\end{document}
